\scanpage{195}

il existe un \inserted{voisinage} ouvert \deleted{\unsure{de}} $U$ \deleted{$\ne\varnothing$ de $S$} de $\eta$ dans $S$, tel que pour tt $s\in U$, \inserted{$k(s)$ de car. $p>0$,} \deleted{$E_s$ sur $X_s$, avec la connexion relativement à $k(s)$} $p$-courbure de $E_s$ sur $X_s/k(s)$ soit nulle. \inserted{ou encore, intégrables ---} (Cette dernière condition est d'ailleurs équivalente à la condition analogue, \deleted{où} en se limitant aux pts \uline{fermés} de $U$, mais peu importe].

\uline{Question} La \uuline{réciproque est-elle vraie} ?

\uline{Cas \deleted{\unsure{des}} où} $E$ \uline{est de rang} $1$. Notons que \deleted{\unsure{si}} $E_{\eta}$ intégrable signifie que $E$ est associé à un revêtement \unsure{Kummerien|Kummérien} de $X_{\eta}$, i.e. \deleted{obtenu en un module $L$} sa connexion est associé à une trivialisation de $E^{\otimes n}$, pour $n\geq1$ convenable.

La question précédente se décompose en deux :

% vertical rule at left of a)

a) Si la connexion des $E_s$ \inserted{($s\in S$, $s$ de car $>0$)} est intégrable, est-il vrai que l'image de $E$ dans $\mathrm{Pic}(X_\eta)$ est de torsion ?

\deleted{b)} Si la réponse est affirmative, on est ramené dans le Pb au cas où $E\simeq\underline O_X$ (comme ne uniquement, pas la connexion, à priori).

% vertical rule at left of b)

b) Si $\omega$ est une $1$-forme différentielle \inserted{[fermée]} sur $X/S$, et si les $\omega_s$ sur les $X_s/s$ ($s\in S$, $s$ de car $>0$) sont \uline{loc.} des différentielles logarithmiques, alors est-il vrai que \inserted{(un multiple entier de)} $\omega_\eta$ est une diff. \uline{log}arithmique ?

\margin{ce qui se vérifie par l'opération de Cartier $\widetilde C\omega_s=\widetilde\omega_s$}

% archive sheet number: 194
